\input{../tex-inputs/latex-leseansicht-vorspann}

\section[Theodor von Sosnosky an Arthur Schnitzler, 30. 1. 1901]{L04229 Theodor von Sosnosky an Arthur Schnitzler, 30. 1. 1901}
\nopagebreak\mylabel{L04229v}
\rehead{ }\normalsize\beginnumbering\briefempfaengerindex{Schnitzler, Arthur@\textsc{Schnitzler, Arthur}!zzzSosnosky, Theodor von@\emph{von Theodor von Sosnosky}!1901-01-301@{30. 1. 1901}|(be}
\toendnotes[C]{\smallbreak\pagebreak[2]}
\correspDesc{Versand  durch Theodor von Sosnosky am 30. 1. 1901 in Kremsmünster
\newline{}Erhalt  durch Arthur Schnitzler im Zeitraum [31. 1. 1901
                  – 4. 2. 1901?] in Wien}\toendnotes[C]{\smallbreak}
\Standort{DLA, A:Schnitzler, HS.1985.1.4640.}
\physDesc{Brief, 1 Blatt, 3 Seiten, 2135 Zeichen
\newline{}Handschrift: blaue Tinte, deutsche Kurrent
\newline{}Schnitzler: mit rotem Buntstift eine Unterstreichung }\toendnotes[C]{\smallbreak}
\pstart
           \raggedleft{}{\pb}dzt \textsc{Kremsmünster\oindex{Kremsmünster@\textbf{Kremsmünster}, \emph{Verwaltungsgebiet}|pw}}, 30/I. 1901\pend
           
\pstart{}Sehr geehrter Herr Doctor!\pend\vspace{0.5em}
\pstart
           Anbei erlaube ich mir Ihnen ein Exemplar meiner Anthologie\pwindex{deutsche Lyrik des 19. Jahrhunderts. Eine poetische Revue@\emph{Die deutsche Lyrik des 19. Jahrhunderts. Eine poetische Revue}|pwv} zu überſenden, dem ich durch eine kleine
               Inſchrift einen perſönlichern Charakter zu geben ſo frei war.\pend
           
\pstart
           Es iſt dies meine urſprüngliche Abſicht geweſen und das für Sie beſtimmte Buch\pwindex{deutsche Lyrik des 19. Jahrhunderts. Eine poetische Revue@\emph{Die deutsche Lyrik des 19. Jahrhunderts. Eine poetische Revue}|pwv} war ſchon eingepackt und
               adreßiert als \label{K_L04229-1v}\edtext{\textsc{Cottas\orgindex{J.G. Cotta’sche Buchhandlung Nachfolger@J.G. Cotta’sche Buchhandlung Nachfolger|pw}} Brief}{\lemma{\textnormal{\emph{Cottas Brief}}}\Cendnote{\textnormal{nicht überliefert}}}\label{K_L04229-1} kam,
               worin er die Abſendung an alle Autoren ankündigte, worauf ich meine Abſicht aufgab
               und das Buch\pwindex{deutsche Lyrik des 19. Jahrhunderts. Eine poetische Revue@\emph{Die deutsche Lyrik des 19. Jahrhunderts. Eine poetische Revue}|pwv} anderweitig
               verwendete. Da ich aber Ihrem heute eingetroffenen liebenswürdigen
                  \label{K_L04229-2v}\edtext{Briefe}{\lemma{\textnormal{\emph{Briefe}}}\Cendnote{\textnormal{XXXX Auszeichnungsfehler: Dokument L04237 nicht gefunden.}}}\label{K_L04229-2} entnehme, daß \textsc{Cotta\orgindex{J.G. Cotta’sche Buchhandlung Nachfolger@J.G. Cotta’sche Buchhandlung Nachfolger|pw}} noch nichts geſandt hat, ſo ſende ich Ihnen perſönlich ein Buch\pwindex{deutsche Lyrik des 19. Jahrhunderts. Eine poetische Revue@\emph{Die deutsche Lyrik des 19. Jahrhunderts. Eine poetische Revue}|pwv}; zugleich {\pb}verſtändigte ich \textsc{Cotta\orgindex{J.G. Cotta’sche Buchhandlung Nachfolger@J.G. Cotta’sche Buchhandlung Nachfolger|pw}} davon, damit Sie nicht durch ein Duplikat\pwindex{deutsche Lyrik des 19. Jahrhunderts. Eine poetische Revue@\emph{Die deutsche Lyrik des 19. Jahrhunderts. Eine poetische Revue}|pwv} beläſtigt werden.\pend
           
\pstart
           Daß Ihnen mein \label{K_L04229-3v}\edtext{\textsc{\textsc{Essai}\pwindex{Sosnosky, Theodor von 4.\,1.\,1866 Budapest – 13.\,2.\,1943 Wien@\textsc{Sosnosky, Theodor von} (4.\,1.\,1866 Budapest – 13.\,2.\,1943 Wien), \emph{Schriftsteller}!Jung-Wien«. Studienblätter@\strich\emph{»Jung-Wien«. Studienblätter}|pwv}} in der Nord. Allgem.\pwindex{Norddeutsche Allgemeine Zeitung@\emph{Norddeutsche Allgemeine Zeitung}|pw}}{\lemma{\textnormal{\emph{Essai … Allgem.}}}\Cendnote{\textnormal{\emph{»Jung-Wien«. Studienblätter von Theodor v.
                        Sosnosky}\pwindex{Sosnosky, Theodor von 4.\,1.\,1866 Budapest – 13.\,2.\,1943 Wien@\textsc{Sosnosky, Theodor von} (4.\,1.\,1866 Budapest – 13.\,2.\,1943 Wien), \emph{Schriftsteller}!Jung-Wien«. Studienblätter@\strich\emph{»Jung-Wien«. Studienblätter}|pwk}. In: Beilage zur \emph{Norddeutschen Allgemeinen Zeitung}\pwindex{Norddeutsche Allgemeine Zeitung@\emph{Norddeutsche Allgemeine Zeitung}|pwk}, Jg. 40, Nr. 17a,
                        20. 1. 1901, S. [1]–2, Nr. 20, 26. 1. 1901, S. [1]
                     und Nr. 29a, 3. 2. 1901, S. [1]–2.}}}\label{K_L04229-3} zu Geſichte
                  geko{\geminationm}en, hat mich nicht wenig überraſcht, da das Blatt\pwindex{Norddeutsche Allgemeine Zeitung@\emph{Norddeutsche Allgemeine Zeitung}|pwv} in den Wiener\oindex{Wien@\textbf{Wien}, \emph{Verwaltungsgebiet}|pw}{ }\textsc{Cafés} faſt nirgend aufliegt; und Sie haben vom Artikel\pwindex{Sosnosky, Theodor von 4.\,1.\,1866 Budapest – 13.\,2.\,1943 Wien@\textsc{Sosnosky, Theodor von} (4.\,1.\,1866 Budapest – 13.\,2.\,1943 Wien), \emph{Schriftsteller}!Jung-Wien«. Studienblätter@\strich\emph{»Jung-Wien«. Studienblätter}|pwv} ſogar mehr geleſen als
               ich ſelber, der ich bisher nur den \uline{1.} Teil\pwindex{Sosnosky, Theodor von 4.\,1.\,1866 Budapest – 13.\,2.\,1943 Wien@\textsc{Sosnosky, Theodor von} (4.\,1.\,1866 Budapest – 13.\,2.\,1943 Wien), \emph{Schriftsteller}!Jung-Wien«. Studienblätter@\strich\emph{»Jung-Wien«. Studienblätter}|pwv} erhalten habe.\pend
           
\pstart
           Ich erhielt ihn – was faſt romanhaft zufällig ausſieht – an demſelben Tage, an dem
               ich Ihnen darüber geſchrieben hatte.\pend
           
\pstart
           Wenn er vollſtändig erſchienen iſt – es dürften \label{K_L04229-4v}\edtext{4 Teile}{\lemma{\textnormal{\emph{4 Teile}}}\Cendnote{\textnormal{Der Text\pwindex{Sosnosky, Theodor von 4.\,1.\,1866 Budapest – 13.\,2.\,1943 Wien@\textsc{Sosnosky, Theodor von} (4.\,1.\,1866 Budapest – 13.\,2.\,1943 Wien), \emph{Schriftsteller}!Jung-Wien«. Studienblätter@\strich\emph{»Jung-Wien«. Studienblätter}|pwkv}
                  wurde in drei Folgen geliefert.}}}\label{K_L04229-4} werden – will ich mir erlauben, ihn Ihnen
               als Ganzes \label{K_L04229-5v}\edtext{zu überſenden}{\lemma{\textnormal{\emph{zu übersenden}}}\Cendnote{\textnormal{Ein Exemplar des Textes\pwindex{Sosnosky, Theodor von 4.\,1.\,1866 Budapest – 13.\,2.\,1943 Wien@\textsc{Sosnosky, Theodor von} (4.\,1.\,1866 Budapest – 13.\,2.\,1943 Wien), \emph{Schriftsteller}!Jung-Wien«. Studienblätter@\strich\emph{»Jung-Wien«. Studienblätter}|pwkv} befindet sich in
                     Schnitzlers Sammlung von
                  Zeitungsausschnitten: Arthur Schnitzler: \emph{Archiv der Zeitungsausschnitte}, \url{https://schnitzler\_zeitungen.acdh.oeaw.ac.at/224251\_0001}.}}}\label{K_L04229-5} mit der Bitte ihn zu behalten.\pend
           
\pstart
           Sie dürfen, geehrter Herr Doktor, verſichert ſein, daß mein Urteil darin ganz ehrlich
               meine Überzeugung \substVorne{}\textsuperscript{iſt}\substDazwischen{}ausdrückt\substHinten{}, und ich glaube, daß Ihr \strikeout{beſcheidener}{ }{\pb}Einwand, es sei in Manchem \uline{zu} günſtig
               ausgefallen, zu beſcheiden ist; hab’ ich doch meine Bedenken, die allerdings nur ſehr
               vereinzelt ſind, unumwunden, fast derb geäußert.\pend
           
\pstart
           Ich bitte auch \label{K_L04229-6v}\edtext{eine Äußerung}{\lemma{\textnormal{\emph{eine Äußerung}}}\Cendnote{\textnormal{Zum Aspekt, ob die Autoren der Gruppierung
                     \emph{Jung Wien}\orgindex{Jung Wien@Jung Wien|pwk} »einen \so{echten} Klang jeder wienerische\oindex{Wien@\textbf{Wien}, \emph{Verwaltungsgebiet}|pw} Poesie« träfen, schreibt Sosnosky: »Bei Schnitzler{ }\so{glaubt} man sie wohl zu hören, aber es ist doch nur
                     eine, allerdings recht glückliche Imitation, so etwa, als ob eine gute
                     Zigeunerkapelle einen Wiener\oindex{Wien@\textbf{Wien}, \emph{Verwaltungsgebiet}|pw} Walzer spielte;
                     ein feineres Ohr muß den fremden Grundklang heraushören.« \emph{»Jung-Wien«. Studienblätter von Theodor v.
                        Sosnosky}\pwindex{Sosnosky, Theodor von 4.\,1.\,1866 Budapest – 13.\,2.\,1943 Wien@\textsc{Sosnosky, Theodor von} (4.\,1.\,1866 Budapest – 13.\,2.\,1943 Wien), \emph{Schriftsteller}!Jung-Wien«. Studienblätter@\strich\emph{»Jung-Wien«. Studienblätter}|pwk}. In: Beilage zur \emph{Norddeutschen Allgemeinen Zeitung}\pwindex{Norddeutsche Allgemeine Zeitung@\emph{Norddeutsche Allgemeine Zeitung}|pwk}, Jg. 40, Nr. 29a,
                        3. 2. 1901, S. 2. Auf der vorangegangenen Seite hatte er
                  den Autoren bereits »Mangel an Natürlichkeit«, »Mangel an Humor«, »eine gewisse
                  Blasirtheit, Müdheit und Freudlosigkeit«, »eine Art Licht- und Luftscheue« sowie
                  »Mangel an Kraft« bescheinigt.}}}\label{K_L04229-6}, die ich am Schluße des ganzen \textsc{\textsc{Essais}\pwindex{Sosnosky, Theodor von 4.\,1.\,1866 Budapest – 13.\,2.\,1943 Wien@\textsc{Sosnosky, Theodor von} (4.\,1.\,1866 Budapest – 13.\,2.\,1943 Wien), \emph{Schriftsteller}!Jung-Wien«. Studienblätter@\strich\emph{»Jung-Wien«. Studienblätter}|pwv}} gemacht habe, lediglich als \uline{allgemein\strikeout{}}\strikeout{e Anſicht}, nicht als \uline{perſönlich} aufzufaſſen. Nichts wäre mir peinlicher, als Sie zu verletzen.
                  (Schwarzkopf\pwindex{Schwarzkopf, Gustav 7.\,11.\,1853 Wien – 13.\,11.\,1939 ebd.@\textsc{Schwarzkopf, Gustav} (7.\,11.\,1853 Wien – 13.\,11.\,1939 ebd.), \emph{Schriftsteller}|pw} kennt meine diesbezügl.
               Anſichten ſo ziemlich).\pend
           
\pstart
           Der beſte Gradmeſſer für meine Hochschätzung Ihrer Kunst ist übrigens, wie mir
               scheint, nicht ſo ſehr meine Anerkennung in Worten als vielmehr die Folie, die Sie
               durch die übrigen beſprochenen \label{K_L04229-7v}\edtext{Autoren\pwindex{Bahr, Hermann 19.\,7.\,1863 Linz – 15.\,1.\,1934 München@\textsc{Bahr, Hermann} (19.\,7.\,1863 Linz – 15.\,1.\,1934 München), \emph{Schriftsteller, Kritiker}|pwv}\pwindex{Hofmannsthal, Hugo von 1.\,2.\,1874 Wien – 15.\,7.\,1929 Rodaun@\textsc{Hofmannsthal, Hugo von} (1.\,2.\,1874 Wien – 15.\,7.\,1929 Rodaun), \emph{Schriftsteller}|pwv}\pwindex{Altenberg, Peter 9.\,3.\,1859 Wien – 8.\,1.\,1919 ebd.@\textsc{Altenberg, Peter} (9.\,3.\,1859 Wien – 8.\,1.\,1919 ebd.), \emph{Schriftsteller}|pwv}\pwindex{Dörmann, Felix 29.\,5.\,1870 Wien – 26.\,10.\,1928 ebd.@\textsc{Dörmann, Felix} (29.\,5.\,1870 Wien – 26.\,10.\,1928 ebd.), \emph{Schriftsteller}|pwv}\pwindex{Salten, Felix 6.\,9.\,1869 Budapest – 8.\,10.\,1945 Zürich@\textsc{Salten, Felix} (6.\,9.\,1869 Budapest – 8.\,10.\,1945 Zürich), \emph{Schriftsteller, Journalist, Chefredakteur}|pwv}\pwindex{Wolff, Ludwig 7.\,3.\,1876 Bielsko-Biała – nach 1958 Vereinigte Staaten von Amerika [USA]@\textsc{Wolff, Ludwig} (7.\,3.\,1876 Bielsko-Biała – nach 1958 Vereinigte Staaten von Amerika [USA]), \emph{Schriftsteller, Dramaturg}|pwv}}{\lemma{\textnormal{\emph{Autoren}}}\Cendnote{\textnormal{Der \emph{Essay}\pwindex{Sosnosky, Theodor von 4.\,1.\,1866 Budapest – 13.\,2.\,1943 Wien@\textsc{Sosnosky, Theodor von} (4.\,1.\,1866 Budapest – 13.\,2.\,1943 Wien), \emph{Schriftsteller}!Jung-Wien«. Studienblätter@\strich\emph{»Jung-Wien«. Studienblätter}|pwk} behandelte die Schriftsteller Hermann Bahr\pwindex{Bahr, Hermann 19.\,7.\,1863 Linz – 15.\,1.\,1934 München@\textsc{Bahr, Hermann} (19.\,7.\,1863 Linz – 15.\,1.\,1934 München), \emph{Schriftsteller, Kritiker}|pwk}, Arthur Schnitzler,
                     Hugo von Hofmannsthal\pwindex{Hofmannsthal, Hugo von 1.\,2.\,1874 Wien – 15.\,7.\,1929 Rodaun@\textsc{Hofmannsthal, Hugo von} (1.\,2.\,1874 Wien – 15.\,7.\,1929 Rodaun), \emph{Schriftsteller}|pwk}, Richard Beer-Hofmann\pwindex{Beer-Hofmann, Richard 11.\,7.\,1866 Wien – 26.\,9.\,1945 New York City@\textsc{Beer-Hofmann, Richard} (11.\,7.\,1866 Wien – 26.\,9.\,1945 New York City), \emph{Schriftsteller}|pwk}, Peter Altenberg\pwindex{Altenberg, Peter 9.\,3.\,1859 Wien – 8.\,1.\,1919 ebd.@\textsc{Altenberg, Peter} (9.\,3.\,1859 Wien – 8.\,1.\,1919 ebd.), \emph{Schriftsteller}|pwk}, Felix Dörmann\pwindex{Dörmann, Felix 29.\,5.\,1870 Wien – 26.\,10.\,1928 ebd.@\textsc{Dörmann, Felix} (29.\,5.\,1870 Wien – 26.\,10.\,1928 ebd.), \emph{Schriftsteller}|pwk}, Felix Salten\pwindex{Salten, Felix 6.\,9.\,1869 Budapest – 8.\,10.\,1945 Zürich@\textsc{Salten, Felix} (6.\,9.\,1869 Budapest – 8.\,10.\,1945 Zürich), \emph{Schriftsteller, Journalist, Chefredakteur}|pwk} und Ludwig Wolff\pwindex{Wolff, Ludwig 7.\,3.\,1876 Bielsko-Biała – nach 1958 Vereinigte Staaten von Amerika [USA]@\textsc{Wolff, Ludwig} (7.\,3.\,1876 Bielsko-Biała – nach 1958 Vereinigte Staaten von Amerika [USA]), \emph{Schriftsteller, Dramaturg}|pwk}.}}}\label{K_L04229-7} erhalten, die ich nichts weniger als
               zart angefaßt und von denen mir nur einer wert ist – \textsc{Beer-Hofmann\pwindex{Beer-Hofmann, Richard 11.\,7.\,1866 Wien – 26.\,9.\,1945 New York City@\textsc{Beer-Hofmann, Richard} (11.\,7.\,1866 Wien – 26.\,9.\,1945 New York City), \emph{Schriftsteller}|pw}}.\pend
           
\pstart
           Mit dem Ausdrucke vorzüglicher{\\[\baselineskip]}Hochachtung{\\[\baselineskip]}\spacefill\mbox{Theodor vSonosky}\pend
           \leftskip=0em{}\selectlanguage{ngerman}\endnumbering\briefempfaengerindex{Schnitzler, Arthur@\textsc{Schnitzler, Arthur}!zzzSosnosky, Theodor von@\emph{von Theodor von Sosnosky}!1901-01-301@{30. 1. 1901}|)be}\mylabel{L04229h}
\begin{anhang}
\end{anhang}\newcommand{\dateiname}{L04229}\newcommand{\titel}{Theodor von Sosnosky an Arthur Schnitzler, 30. 1. 1901}\newcommand{\editorInnen}{Herausgegeben von Jahnke, SelmaMüller, Martin Anton}\input{../tex-inputs/latex-leseansicht-abspann}
